\section{Additional Specification}
\label{app:additional-specification}

\subsection{Task Output Format}
\label{app:task-output-format}

In \name{}, we define two possible task-instruction types, namely \(\mathcal{I}_{\mathrm{direct}}\) and \(\mathcal{I}_{\mathrm{code}}\), with two corresponding evaluation methods.
In the case of \(\mathcal{I}_{\mathrm{direct}}\), we instruct the agent to return the correct predictions directly; namely, the agent writes a file \texttt{results.json} of the form
\[
\begin{array}{l}
\{ \\
\quad \texttt{"<filename>.parquet"}: [ \\
\quad\quad \texttt{"<top-1 label>"}, \\
\quad\quad \texttt{"<optional lower-confidence label>"} \\
\quad ] \\
\}.
\end{array}
\]
We use this solution to evaluate test samples.
In the case of \(\mathcal{I}_{\mathrm{code}}\), the agent writes a self-contained Python file \texttt{rule.py} defining
\[
\texttt{predict(df) -> list[str]}.
\]

The input \texttt{df} is a single sample represented as a pandas DataFrame.
The returned list follows the same ranked-label convention: the first element is the top-1 prediction, and any later elements are lower-confidence alternatives.
All returned labels must exactly match one of the allowed agent-facing label strings for the corresponding environment.

\subsection{Agentic Prompt Template}
\label{app:agentic-prompt}

The agentic prompt template is provided in the Prompt Example Combinations appendix.
